\originalchapter{来源对应与原件校正}
\originalsection{讲义对应}
所有下列讲义均为同一课程的实际 Typed Notes, 文件名为 \texttt{MIT8\_03SCF16\_LecN.pdf}. 本册按物理内容重组, 不是每页逐字镜像. 课堂演示与重复复习合并, 主结论、适用条件和有实质教学作用的例子保留并展开.
\begin{longtable}{p{.12\linewidth}p{.39\linewidth}p{.36\linewidth}}
\toprule 原讲次 & 原材料实质主题 & 本册定位\\\midrule\endhead
1 & 受力、初值、势能线性化、复表示 & 第\ref{ch:osc}章\\
2--3 & LC、物理摆、阻尼、驱动与共振 & 第\ref{ch:osc}--\ref{ch:drive}章\\
4--6 & 正常模、非等质量、拍、耦合驱动 & 第\ref{ch:modes}章\\
7--8 & 反射和平移对称、离散链、有限边界、连续极限 & 第\ref{ch:chain}章\\
9--10 & 驻波、傅里叶级数、行波、能量与连接 & 第\ref{ch:string}章\\
11 & 纵波、绝热声速、管内边界与驱动 & 第\ref{ch:sound}章\\
12 & Maxwell波、矢量场、平面波、能流 & 第\ref{ch:em}章\\
13--14 & 刚性弦、群相速度、截止、Fourier与调幅 & 第\ref{ch:disp}--\ref{ch:fourier}章\\
15--16 & 方差不确定关系、二维网、膜与三维波、Snell & 第\ref{ch:fourier}--\ref{ch:interface}章\\
17--18 & 线圆椭圆偏振、非偏振、投影、波片、加速辐射 & 第\ref{ch:em}、\ref{ch:polar}章\\
19 & 极化磁化、介质、界面系数、Brewster角 & 第\ref{ch:em}、\ref{ch:interface}章\\
20--22 & 薄膜、双缝、多源雷达、单缝、圆孔与分辨率 & 第\ref{ch:diffraction}章\\
23 & 物质波、路径信息、势阱与量子色散 & 第\ref{ch:quantum}章\\
24 & 复习视频, 未公开对应文字讲义 & 仅列课历, 不冒称已翻译录像\\\bottomrule
\end{longtable}

\originalsection{作业与练习卷}
作业1--10依原编号完整列于作业章, 共44道大题. 实际Exam1、Exam2共7道大题另列完整章, 全部原小问与AI独立解答纳入, 没有把练习卷冒充实际考试. 原PS8仅三题; PS9.2的小问是a,b,d,e, 没有c, 本稿不新增; PS10末题原印Problem 5, 本册用10.5索引并保留原印说明. 作业中的参考讨论不替代具体题面. 练习卷共3、6、6、6、3题, 分别为期中练习一、二与期末练习一、二、三, 总计24道; 五份答案均来自官方 \href{https://ocw.mit.edu/courses/8-03sc-physics-iii-vibrations-and-waves-fall-2016/resources/exam-solutions/}{Exam Solutions} 页面. 资源索引中某些“not available”标签与实际公开PDF不一致, 本册依据实际原件. 实际考试题面来自官方总下载页的独立Exam1/Exam2资源, 二者封面均写Fall 2016; 五份练习答案不作为它们的官方答案. 原公式表支持正文常用公式, 不作为额外未见题面来源. 本册所收作业与试卷共75道大题. 每份原件的准确URL、大小、SHA-256及下载校验见公开清单 \texttt{manifests/waves-and-optics-sources.json}.

\originalsection{校正与归属}
中文推导使用量纲、边界与初值复核, 未机械复制下列原件错式: 第3讲响应分母漏写 $(\omega_0^2-\omega_d^2)$ 的平方, 高频相位应取 $\pi$ 而非由正切主值决定; 第7讲“任意双分量系统”共用对称模的表述须限定对称结构; 第8讲有限固定链有 $N$ 个非零模; 第11讲压力符号及受驱管的共振零点需与所写边界一致; 第14讲delta归一化含 $2\pi$; 第15讲载波以相速度而非群速度前进; 第16讲首式色散漏平方根及系数, 第18讲非相对论辐射总功率系数应为 $1/(6\pi\epsilon_0c^3)$; 第19讲Brewster角的单位透射是功率而非电场振幅; 圆孔强度的中心归一需 $2J_1(\beta)/\beta$.

练习卷的校正随题解注明, 保留原官方解答出处, 不把本稿纠错称为官方修订. PS1.4的原印势能量纲与括号问题按题面原样讨论, 不默改成另一题. Practice Final1的封面内部标题为Exam1且第5题分值不一致, 按门户识别卷名并记录差别. Practice Final3答案扫描中的Spring 2014只说明该原件的标注, 不自动套给另外四卷.

\begin{thebibliography}{9}
\bibitem{ocw} Yen-Jie Lee. \emph{8.03SC Physics III: Vibrations and Waves}. MIT OpenCourseWare, Fall 2016. \url{https://ocw.mit.edu/courses/8-03sc-physics-iii-vibrations-and-waves-fall-2016/}.
\bibitem{georgi} Howard Georgi. \emph{The Physics of Waves}. Benjamin Cummings, 1992. 官方课程指定教材, 本轮不以整书为翻译范围.
\bibitem{faq} MIT OpenCourseWare. \emph{How is “all rights reserved” content different from the rest of OCW content?} 2025年6月23日更新. \url{https://mitocw.zendesk.com/hc/en-us/articles/4414756181403-How-is-all-rights-reserved-content-different-from-the-rest-of-OCW-content}.
\end{thebibliography}
